\chapter{Forschungsstand}
\label{ch:literature}

Dieses Kapitel ordnet die vorliegende Arbeit in den bisherigen Forschungsstand
ein. Es beginnt mit den deutschen Untersuchungen zur Festigkeit von
\gls{altholz}, wendet sich anschliessend internationalen Arbeiten zu,
betrachtet die verfuegbaren zerstoerungsfreien Pruefverfahren und ihre
Ortsaufloesung und benennt zuletzt die Forschungsluecke, an der die eigene
Untersuchung ansetzt.

\section{Deutsche Untersuchungen zur Altholzfestigkeit}
\label{sec:lit-deutsch}

Die deutschsprachige Forschung hat sich ab Mitte der 1980er Jahre systematisch
mit der Frage befasst, ob altes Konstruktionsholz eine geringere Festigkeit
aufweist als neues Holz derselben Holzart. Die zentrale Referenz ist die Arbeit
von \textcite{rug1989}. Untersuchungen an fehlerfreien Kleinproben und an
Laborproben in Schnittholzabmessungen fuehrten dort zu dem Ergebnis, dass
zwischen unbeschaedigtem Altholz und Neuholz derselben Holzart kein
Festigkeitsunterschied besteht. Diese Aussage ist bis heute die Grundlage der
meisten praktischen Argumentationen fuer eine tragende Wiederverwendung: Nicht
das Alter des Holzes begrenzt die \gls{restfestigkeit}, sondern der konkrete
Schaden am Bauteil. Anzumerken ist, dass die Literaturrecherche zu dieser
Arbeit sich fuer diesen Befund auf eine Sekundaerdarstellung stuetzt und der
Volltext nicht geprueft werden konnte.

Parallel dazu entstand in Karlsruhe im Rahmen des Sonderforschungsbereichs 315
eine Reihe von Arbeiten von Ehlbeck und Goerlacher. Die ersten Ergebnisse von
Festigkeitsuntersuchungen an altem Konstruktionsholz \parencite{ehlbeck1988}
und die spaetere Auseinandersetzung mit der Problematik bei der Beurteilung der
Tragfaehigkeit von altem Konstruktionsholz \parencite{ehlbeck1990} verschieben
den Schwerpunkt von der Materialfrage zur Bewertungsfrage. Nicht das Holz ist
das Problem, sondern die Unsicherheit darueber, welche Festigkeit einem
eingebauten, nicht zerstoerend pruefbaren Bauteil zugewiesen werden darf. In
denselben Zusammenhang gehoert eine weitere Karlsruher Veroeffentlichung zu
Bohrwiderstandsmessungen an eingebautem Konstruktionsholz. Nach den
vorliegenden Angaben wandten die Autoren dieses Verfahren an eingebauten
Hoelzern an, um Schaeden zu lokalisieren und eine Festigkeitsreserve
abzuschaetzen.

Beide Arbeitsstraenge zusammen ergeben ein bemerkenswertes Bild. Der Befund,
dass unbeschaedigtes Altholz nicht schwaecher ist als Neuholz, ist mit der
\gls{sortierung} nach \textcite{din4074} und der Klassifizierung nach
\textcite{en338} vereinbar, weil beide Regelwerke ohnehin nicht das Alter,
sondern sichtbare Merkmale und die \gls{rohdichte} bewerten. Genau deshalb
bleibt aber offen, was die \gls{lastgeschichte} eines Balkens mit dem Material
gemacht hat, denn sie ist am Bauteil nicht sichtbar und in keinem
Sortierkriterium abgebildet.

\section{Internationale Arbeiten}
\label{sec:lit-international}

Die aussagekraeftigste internationale Referenz ist eine estnische Fallstudie an
rund 120 Jahre altem Bauholz aus einer historischen Bibliothek in Tartu
\parencite{tartu}. Methodisch verbindet sie zerstoerungsfreie Voruntersuchungen
mit zerstoerenden Vierpunkt-Biegeversuchen und stellt die Bewertung nach alten
skandinavischen Sortierregeln, nach nordischen Normen (SS 230120, INSTA 142)
und nach italienischen Regeln fuer die Altholzbewertung einander gegenueber.
Die Studie ist damit fuer die vorliegende Arbeit in zwei Hinsichten
interessant: Sie zeigt, dass jahrzehntealtes Bauholz einer regulaeren
Festigkeitspruefung zugaenglich ist, und sie macht die Abhaengigkeit der
Bewertung vom gewaehlten Regelwerk zum ausdruecklichen Gegenstand. Eine dort
berichtete Beobachtung zur guenstigsten Holzgewebequalitaet bei einem Baumalter
von etwa 80 Jahren betrifft das Schlagalter des Baumes und nicht die Alterung
des Holzes im eingebauten Zustand; sie wird hier nur der Vollstaendigkeit wegen
erwaehnt und nicht als Alterungsergebnis gedeutet.

Aus dem nordeuropaeischen Raum liegen ausserdem die Sortierregelwerke selbst
vor, deren experimentelle Absicherung an Altholz nach dem Stand der gesichteten
Literatur duenn bleibt. Eine Arbeit zum Schubmodul von altem Bauholz
\parencite{iforest} zeigt zudem, dass neben Biegefestigkeit und
\acrshort{moe} auch weitere Kennwerte alter Hoelzer Gegenstand aktueller
Untersuchungen sind.

Die US-amerikanische Forschung ist fuer die Fragestellung vor allem als
Grundlagenliteratur einschlaegig. Die Zusammenstellung der Holzeigenschaften im
Handbuch des USDA Forest Service \parencite{usdaTimberBridge} beschreibt die
richtungsabhaengigen Festigkeiten und damit den Rahmen, in dem
\gls{druckzone} und \gls{zugzone} eines Biegetraegers ueberhaupt
unterschiedlich beansprucht werden. Hinzu kommen Arbeiten zum Langzeitverhalten:
Genannt werden Kriechversuche und Modelle der Schadensansammlung, nach denen
ein Holzbauteil kurzzeitig hohe Lasten tragen kann, unter dauernder Last aber
an Tragfaehigkeit verliert (\acrshort{dol}-Effekt). Diese Angaben konnten nicht
am Original geprueft werden und werden hier daher nur als Hinweis auf einen
plausiblen Wirkmechanismus verstanden, nicht als quantitative Grundlage. Nach
denselben Sekundaerangaben umfassen die verfuegbaren Laborstudien zum
\gls{kriechen} ueberdies Zeitraeume von etwa einem Jahr, waehrend die hier
untersuchten Balken jahrzehntelang unter Last standen.

\section{Zerstoerungsfreie Pruefverfahren und ihre Ortsaufloesung}
\label{sec:lit-ndt}

Fuer die Beurteilung eingebauter Hoelzer haben sich zerstoerungsfreie Verfahren
(\acrshort{ndt}) etabliert. Zum Kern gehoeren die Sichtpruefung, die Bestimmung
von Holzfeuchte und Holzart, die Bohrwiderstandsmessung, Laufzeitmessungen
elastischer Wellen (Stresswave und Ultraschall) sowie radarbasierte und
bildgebende Verfahren. Eine internationale Empfehlung fasst diese Methoden fuer
die Bewertung im Bestand zusammen; sie liegt in den gesichteten Quellen ohne
Autorennamen vor und wird daher hier ohne solche genannt. Die Holzfeuchte, die
jede dieser Messungen beeinflusst, wird nach \textcite{din13183} bestimmt.

Entscheidend fuer die vorliegende Arbeit ist nicht, ob diese Verfahren
funktionieren, sondern was sie raeumlich aufloesen. Bohrwiderstands- und
Laufzeitmessungen werden in der Bestandsbeurteilung in der Regel zu einem
Gesamtkennwert je Messstelle verdichtet und nicht nach Hoehenlage im
Querschnitt getrennt ausgewertet. Eine getrennte Aussage fuer den oberen und
den unteren Randbereich desselben Balkens ist damit nur begrenzt moeglich.
Genau daran scheitert bislang eine zonenweise Beurteilung: Selbst ein
sorgfaeltig zerstoerungsfrei geprueftes Bauteil liefert einen Kennwert, der
einen moeglichen Unterschied zwischen frueherer Druck- und frueherer Zugzone
verdeckt. Dass neuere bildgebende Verfahren wie Radar und Ultraschall eine
bessere raeumliche Aufloesung erreichen, legen die gesichteten Quellen nahe,
ohne dass dies hier am Original geprueft werden konnte; belastbar bleibt
vorlaeufig nur der Befund, dass bildgebende Messung und zerstoerende Pruefung
an derselben Stelle selten kombiniert werden. Auf der Seite der zerstoerenden
Pruefung liegt mit dem Vierpunkt-Biegeversuch ein etabliertes Verfahren vor,
fuer das auch neuere Protokolle zur Festigkeitssortierung vorgeschlagen werden
\parencite{fourpoint}; die zugehoerigen Randbedingungen regeln
\textcite{en408} und \textcite{en384}.

\section{Forschungsluecke und Einordnung der vorliegenden Arbeit}
\label{sec:lit-luecke}

Mechanisch ist der Ausgangspunkt gut beschrieben. Holz versagt unter Zug
sproede, weil die Festigkeit der Cellulose-Mikrofibrillen kaum plastische
Reserven hat, waehrend es in der Druckzone duktil reagiert und die Zellen
allmaehlich stauchen \parencite{doitpoms}. Druck- und Zugzone eines
langzeitbelasteten Balkens haben also nicht nur unterschiedliche Spannungen,
sondern auch unterschiedliche Schaedigungsmechanismen erfahren, getrennt durch
die \gls{neutralefaser}.

Dazu findet sich jedoch keine systematische Untersuchung. Es fehlt Literatur,
die Proben aus verschiedenen Hoehenlagen desselben langzeitbelasteten Balkens
entnimmt und deren Festigkeitseigenschaften vergleicht. Soweit ersichtlich
pruefen bisherige Arbeiten entweder das Bauteil als Ganzes oder entnehmen
Proben ohne Bezug zur ehemaligen Lage im Querschnitt. Ebenso fehlen
Untersuchungen dazu, ob und wie sich die neutrale Faser unter langer
Gebrauchslast verlagert, wie ungleichmaessig sich Schaedigung ueber die
Querschnittshoehe verteilt und wie sich Alterungseffekte des Materials von
Effekten der Lastgeschichte trennen lassen. Hinzu kommt, dass die gaengige
Unterscheidung zwischen beschaedigtem und unbeschaedigtem Altholz nach dem
Stand der gesichteten Literatur unscharf bleibt, obwohl gerade sie den Befund
von Rug und Seemann traegt.

Hier setzt die vorliegende Arbeit an. Sie entnimmt Proben gezielt nach
Hoehenlage aus demselben Balken, prueft sie unter vergleichbaren Bedingungen
und vergleicht die Ergebnisse der ehemaligen Druckzone mit denen der ehemaligen
Zugzone. Damit wird die Hoehenlage im Querschnitt von einer bisher
vernachlaessigten Randbedingung zur eigentlichen Untersuchungsvariablen. Der
Beitrag liegt weniger in einem neuen Pruefverfahren als in der konsequenten
Zuordnung jeder Probe zu ihrer Lage und damit zu ihrer Beanspruchungsgeschichte,
also an genau der Stelle, an der die zerstoerungsfreien Verfahren ihre
Ortsaufloesung verlieren.
